% Generated draft: requires mathematical and source review.
\documentclass[11pt]{article}
\usepackage[margin=25mm]{geometry}
\usepackage{fontspec}
\setmainfont{Latin Modern Roman}
\usepackage{amsmath,amssymb,mathtools,unicode-math}
\setmathfont{Latin Modern Math}
\usepackage{enumitem,needspace,xurl,hyperref}
\setlength{\parindent}{0pt}
\setlength{\parskip}{5pt}
\setlength{\emergencystretch}{3em}
\newcommand{\N}{\mathbb{N}}
\newcommand{\Z}{\mathbb{Z}}
\newcommand{\Q}{\mathbb{Q}}
\newcommand{\R}{\mathbb{R}}
\newcommand{\C}{\mathbb{C}}
\newcommand{\F}{\mathbb{F}}
\newcommand{\A}{\mathbb{A}}
\newcommand{\PP}{\mathbb{P}}
\newcommand{\E}{\mathbb{E}}
\newcommand{\eps}{\varepsilon}
\newcommand{\dd}{\,\mathrm{d}}
\begin{document}
\section*{UnsolvedMath ID: 30002552}
\textbf{Source problem code:} catalogue entry
\subsection{Definitions and mathematical statement}
Let \((D,\mathcal A,\mu)\) denote the domain and reference measure of the estimation problem. These are not specified in the supplied record. For a measurable function \(g\), write
\[
\|g\|_{L^p(\mu)}=\left(\int_D |g|^p\,\dd\mu\right)^{1/p}\quad(1\le p<\infty),\qquad
\|g\|_{L^\infty(\mu)}=\operatorname*{ess\,sup}_{x\in D}|g(x)|.
\]
Write
\[
\mathcal H^\beta(Q;D)=\{f:\|f\|_{C^\beta(D)}\le Q\}
\]
for a H\"older ball of smoothness \(\beta>0\) and radius \(Q>0\). This notation is provisional: the domain, boundary conditions, H\"older norm, and convention at integer smoothness are not supplied. Neither the admissible ranges of \(\beta,Q\) nor the observation model and sample-size parameter are specified.

An adaptive estimator here means an estimator constructed from the observations without knowing \(\beta\) or \(Q\). Simultaneous adaptation asks for one estimator sequence that attains the required performance for all the losses and parameter classes under consideration, rather than selecting a different estimator using the unknown parameters or the requested loss.

\textbf{Research questions (permanent catalogue entry).} In the intended statistical experiment, determine whether such an estimator can satisfy the following requirements:
\begin{enumerate}
\item Attain the appropriate adaptive estimation rates over \(\mathcal H^\beta(Q;D)\), simultaneously for every intended \(L^p\) loss and every admissible unknown pair \((\beta,Q)\).
\item Simultaneously attain the requested sharp \(L^2\) and \(L^\infty\) rates for \(f\) and its derivatives. If the domain is Euclidean and ordinary partial derivatives are intended, write \(D^\alpha f\) for the derivative indexed by a multi-index \(\alpha\), with \(D^0f=f\); the admissible derivative orders and their relation to \(\beta\) must be specified before this requirement is well defined.
\item Extend the intended estimation method to a setting that does not assume coefficientwise bounds on wavelet coefficients.
\item Characterize the settings in which simultaneous adaptation is possible and construct an estimator attaining the corresponding performance guarantees.
\end{enumerate}
The record does not specify the statistical experiment, risk functional, benchmark rates, meaning of sharpness, admissible derivative orders, wavelet assumptions, or the method to be generalized. Its original wording asks for ``clean'' rates whereas its clean statement asks for ``sharp'' rates; these need not express the same requirement. Consequently, these are parameterized research questions, not a fully specified theorem or conjecture. No particular rate, risk criterion, or meaning of sharpness is asserted here.
\subsection{Short English statement}
The problem asks whether one estimator, without knowing the H\"older smoothness or radius, can adapt simultaneously across \(L^p\) losses and derivatives, achieve the requested \(L^2\) and \(L^\infty\) performance, and extend beyond coefficientwise wavelet restrictions. It also asks for a characterization of when this is possible and a construction. Essential model assumptions and performance benchmarks are missing, so a unique precise statistical assertion cannot yet be extracted.
\subsection{Sources}
\begin{itemize}
\item Supplied local problem registry record: permanent UnsolvedMath ID \(30002552\), title \emph{Simultaneous Adaptive Estimation Across Norms and Derivatives}, exact source problem code \texttt{catalogue entry}. The supplied collection is \emph{Oberwolfach Reports Open Problems}; the proposed year is 2014, and no individual proposer is identified.
\item Supplied report citation: \emph{Adaptive Statistical Inference} (2014), Oberwolfach Reports, pp.~721--779, DOI 10.4171/owr/2014/13. Original report link: \url{https://doi.org/10.4171/owr/2014/13}. No specific problem page, speaker, or passage is supplied.
\item Supplied literature source labeled ``Related adaptive estimator'': \url{https://arxiv.org/abs/1312.4087}. The record provides no authors, title, or theorem identifying its precise connection to the questions.
\item Literature assessment dated exactly \texttt{2026-08-21}: the supplied curation labels the problem open, states that later work resolves particular simultaneous-adaptation settings, and reports that no complete solution to all four questions was found. This is supplied dated curation, not an independently verified literature assessment. The background explicitly requests independent verification of time-sensitive status and cited claims.
\item The registry supplies statement status \texttt{corrected\_verified}, with the explanation that mathematical content was preserved while notation, typography, or omitted definitions were repaired for publication. This is a supplied editorial qualification, not independent mathematical verification of this draft.
\end{itemize}
\end{document}
